\newcommand{\bbeproject}[1]{
\section{Wer? Wie? Was? - Die Rahmenbedingungen}
\begin{itemize}
    \item Abgabe in Teams bestehend aus je \textbf{drei Personen}
    \item Abgabe am \textbf{#1} über BYCS-Drive.
    \item Freie Themenwahl (im Rahmen eines Spiels)
    \item Umsetzung in BlueJ basierend auf der Vorlage vom Ressourcen-Laufwerk.
    \item Freie Wahl, ob mit Java, Blockly oder einer Kombination programmiert wird.
    \item \textbf{Tipps zur Inspiration} \begin{itemize}
        \item Mit einer Suchmaschine eurer Wahl: "Greenfoot Project", o.ä. (Greenfoot bietet ähnliche Möglichkeiten, wie die Gameboard-Library)
        \item Bei umfangreichem Kopieren eines bereits bestehenden Projekts aus dem Internet wird das als Unterschleif gewertet!
        \item Dokumentation der Gameboard-Library: \href{https://gameboard.valentin-herrmann.com}{gameboard.valentin-herrmann.com}
    \end{itemize}
    \item \textbf{Arbeit zu Hause:} \begin{itemize}
        \item Sollte eigentlich nicht notwendig sein. Falls doch (wegen Fehlzeiten, Zeitplan verschätzt etc.), sind alle benötigten Programme kostenlos verfügbar.
        \item \textbf{BlueJ} ist für alle Desktop-Betriebssysteme verfübar: \href{https://bluej.org}{bluej.org}.\\$\rightarrow$ Eventuell muss eine Java Runtime installiert werden (eigentlich auf allen Systemen installiert, aber man weiß ja nie... Das merkt man dann aber, Lösung googlen und ggf. von Eltern helfen lassen.) 
        \item Die \textbf{Blockly-Erweiterung für BlueJ} ist auf Github unter Releases (rechte Spalte) verfügbar.\\
        Jede Release-Version hat unter der Beschreibung einen Abschnitt \emph{Assets}, dort kann die Installationsdatei \textbf{aktuell nur für Windows} (.msi) heruntergeladen werden: \href{https://github.com/ValentinHerrmann/BlueJBlocklyExtension/}{github.com/ValentinHerrmann/BlueJBlocklyExtension}\\
        Wenn ihr auf einem anderen Betriebssystem mit Blockly arbeiten müsst, meldet euch bitte bei eurer Lehrkraft.
        \item Wenn ihr nur die \textbf{Gameboard-Library} benötigt (weil ihr direkt mit Java und nicht mit Blockly arbeitet), könnt ihr sie für alle Betriebssysteme hier unter Releases rechte Spalte) herunterladen: \href{https://github.com/ValentinHerrmann/Gameboard}{github.com/ValentinHerrmann/Gameboard}\\
        $\rightarrow$ Unter Assets findet ihr bei jedem Release eine .jar Datei, die ihr auf eurem Computer im BlueJ-Installationsordner (den habt ihr bei der BlueJ Installation ausgewählt) unter \emph{/lib/userlib} einfügt.\\
        Unter Windows ist der Installationsordner normalerweise: \emph{C:/Programme/BlueJ}.
    \end{itemize}
\end{itemize}

\section{Abgaben}
Alle Abgaben über BYCS-Drive als Datei-Freigabe für eure Lehrkraft, es muss nur eine Person pro Team abgeben.
\begin{itemize}
    \item Konzept (ca. eine halbe Seite als PDF, keine Note)\begin{itemize}
        \item Beschreibung des geplanten Projekts
        \item Wird nicht separat benotet, fließt jedoch in die Benotung der Dokumentation ein
        \item Abgabe zu Beginn des Projekts (bevor mit der Programmierung begonnen wird)
    \end{itemize}
    \item Projektordner \begin{itemize}
        \item Abgabe des gesamten Projektordners
        \item Bewertung nach Komplexität, Ausgereiftheit, Schlüssigkeit, Fehleranfälligkeit etc.
    \end{itemize}
    \item Dokumentation als PDF \begin{itemize}
        \item Beschreibung des fertigen Produkts (Features, Steuerung, Bedienungsanleitung etc.)
        \item Klassendiagramm zu allen Klassen mit \href{https://apollon.ase.cit.tum.de/}{apollon.ase.cit.tum.de}
        \item Darstellung der run- und crash-Methoden aller Klassen mit Struktogrammen.\\
        z.B. mit \href{https://janishuser.github.io/Structorizer/struct.html}{janishuser.github.io/Structorizer} oder \href{https://struktog.openpatch.org}{struktog.openpatch.org}
        \item Bericht über Arbeits- und Rollenverteilung in der Projektarbeit
        \item Bewertung nach (formaler) Korrektheit, Vollständigkeit, Umgang mit Hürden während des Entwicklungsprozesses und individuellem Engagement. \textbf{Nicht} danach, ob alles geplante genau so umgesetzt wurde.
    \end{itemize}
\end{itemize}
}